\rhead{SE201: Object-Oriented Analysis \& Design}
\phantomsection 
\section*{Object-Oriented Analysis \& Design (SE201)}
\label{major:SE_1}

\noindent \small{\hyperref[main:scheme]{$\leftarrow$ Back to Scheme}} \vspace{0.5cm}

\noindent \textbf{Credits:} 3 (2-0-2)

\subsection*{Course Content}
\noindent\textbf{Unit 1} \\
\textit{Modeling with UML:} Introduction to the Software Development Life Cycle (SDLC) and the Unified Process. The role of Analysis vs. Design. Unified Modeling Language (UML) 2.0. Functional Modeling: Use Case Diagrams and Scenarios. Structural Modeling: Class Diagrams, Object Diagrams, and Component Diagrams. Behavioral Modeling: Sequence Diagrams, Activity Diagrams, and State Machine Diagrams.

\vspace{0.3cm}

\noindent\textbf{Unit 2} \\
\textit{Design Principles and GRASP:} Fundamental design concepts: Abstraction, Encapsulation, Modularity, Hierarchy. Metrics for Design Quality: Coupling and Cohesion. General Responsibility Assignment Software Patterns (GRASP): Creator, Information Expert, Low Coupling, High Cohesion, and Controller patterns.

\vspace{0.3cm}

\noindent\textbf{Unit 3} \\
\textit{Architectural Principles and SOLID:} The SOLID Principles of Object-Oriented Design: Single Responsibility Principle (SRP), Open/Closed Principle (OCP), Liskov Substitution Principle (LSP), Interface Segregation Principle (ISP), and Dependency Inversion Principle (DIP). Law of Demeter. Dependency Injection and Inversion of Control (IoC).

\vspace{0.3cm}

\noindent\textbf{Unit 4} \\
\textit{GoF Design Patterns (Creational and Structural):} Introduction to Design Patterns. Creational Patterns: Singleton (Thread-safe implementation), Factory Method, Abstract Factory, and Builder. Structural Patterns: Adapter (Class vs. Object adapters), Composite (Tree structures), Decorator (Dynamic responsibilities), Facade, and Proxy.

\vspace{0.3cm}

\noindent\textbf{Unit 5} \\
\textit{GoF Design Patterns (Behavioral) and Refactoring:} Behavioral Patterns: Observer (Publish-Subscribe), Strategy (Algorithm encapsulation), Command, Iterator, and State Pattern. Introduction to Code Smells and Refactoring techniques. Anti-patterns. Mapping Design to Code.

\vspace{0.5cm}

\newpage

\subsection*{Laboratory Content}
\noindent\textbf{Lab 1: Requirement Analysis and Use Cases} \
Focuses on identifying actors and use cases from a problem statement (e.g., Library Management System) and creating a comprehensive Use Case Diagram. \\

\vspace{0.2cm}

\noindent\textbf{Lab 2: Static Structure Modeling} \
Focuses on domain modeling. Students will create detailed Class Diagrams identifying attributes, operations, and relationships (Association, Aggregation, Composition, Generalization). \\

\vspace{0.2cm}

\noindent\textbf{Lab 3: Interaction and Behavior Modeling} \
Focuses on dynamic system behavior. Students will design Sequence Diagrams to model specific scenarios (e.g., "Withdraw Cash") and Activity Diagrams for business workflows. \\

\vspace{0.2cm}

\noindent\textbf{Lab 4: Implementing SOLID Principles} \
Focuses on code quality. Students will be given a "bad" codebase violating SOLID principles and must refactor it to adhere to SRP, OCP, and LSP. \\

\vspace{0.2cm}

\noindent\textbf{Lab 5: Creational Patterns (Factory and Singleton)} \
Focuses on object creation mechanisms. Students will implement a Database Connection pool using Singleton and a UI element generator using the Factory method. \\

\vspace{0.2cm}

\noindent\textbf{Lab 6: Structural Patterns (Adapter and Decorator)} \
Focuses on class composition. Students will implement the Adapter pattern to integrate a legacy third-party library and the Decorator pattern to add features to a text editor (e.g., scrolling, borders). \\

\vspace{0.2cm}

\noindent\textbf{Lab 7: Behavioral Patterns (Strategy and Observer)} \
Focuses on object communication. Students will implement a Payment Processing system using the Strategy pattern and a Stock Market ticker using the Observer pattern. \\

\vspace{0.2cm}

\noindent\textbf{Lab 8: The Composite Pattern} \
Focuses on hierarchical structures. Students will design a File System structure (Folders and Files) where individual objects and compositions are treated uniformly. \\

\vspace{0.2cm}

\noindent\textbf{Lab 9: State Pattern Implementation} \
Focuses on state-dependent behavior. Students will implement the logic for a Vending Machine or an Order Processing System where behavior changes based on internal state. \\

\vspace{0.2cm}

\noindent\textbf{Lab 10: Capstone: Full System Design} \
Focuses on the end-to-end process. Students will select a project (e.g., E-Commerce Backend), produce full UML documentation, and implement the core modules using at least 3 distinct design patterns.

\newpage
\rhead{Curriculum Schema}